\chapter{Nyelvi alapok és HTTP-határok}

\noindent\textbf{Tanulási út: Gyors út: alap.} Tanuld meg a napi munkához szükséges minimális Rust-szerű nyelvi és HTTP surface-t.\par\medskip

\rustbridgeblockhu{Használd tovább a Rustból ismert lokális változókat, \code{Vec<T>}, \code{Option<T>}, \code{String}, skalárokat, metódusokat és \code{match}-et.}{A Velran tulajdona a verified web surface: callable attribútumok, route/input deklarációk, capabilityk, HTML/JSON határok és a biztonságosan lowerelhető részhalmaz.}{A hétköznapi compute-szintaxis nem rivális DSL. Ha egy normál Rust-alakú kifejezés elfogadott, azt használd, ne keress külön Velran-specifikus írásmódot.}


Az A--Z példa már megmutatta a fő építőelemeket. Első olvasáskor koncentrálj a modulokra, typed értékekre, handlerekre, route-okra, opcionális/listaértékekre és biztonságos HTML-re; a részletes aritmetikai és szövegkezelési katalógushoz szükség esetén térj vissza. Munkaszabály: \emph{a bemenetet a határon parse-old és normalizáld, bent typed értékekkel dolgozz, explicit query-vel perzisztálj, és minél később renderelj}.

\section{Forrásstruktúra és statement-határok}

\subsection*{Modulok}

A \code{main.vrn} a modulgráf belépési pontja:

\begin{lstlisting}[language=Velran]
mod models;
mod queries;
mod components;
mod pages;
mod actions;
\end{lstlisting}

A \code{mod} application-root relatív modult tölt be, de a deklarációit nem önti a globális névtérbe. Modulhatáron át az ordinary library-elemek kvalifikált nevet, például \code{models::Article}-t használnak, és alapból privátak; explicit \code{pub} tehet publikus API-vá structot, enumot, pure függvényt, inherent metódust vagy struct-mezőt. A webes authority továbbra is route/capability policyból származik, a \code{pub mod} re-export szemantika pedig külön modul/import kérdés.

\subsection*{Modul-névterek és modulok közötti hívások}

A moduldeklaráció egy application-root relatív forrásmodult vesz fel az explicit source graphba; nem másolja annak deklarációit az aktuális scope-ba. A kanonikus leképezés közvetlen:

\begin{lstlisting}
models.vrn              -> models
catalog/queries.vrn     -> catalog::queries
catalog/pages.vrn       -> catalog::pages
\end{lstlisting}

Egy belépési pont például több modult tölthet be:

\begin{lstlisting}[language=Velran]
mod models;
mod catalog::queries;
mod catalog::pages;

route catalogIndex GET "/catalog"
    public => catalog::pages::index;
\end{lstlisting}

Modulhatáron át kvalifikált nevet használj. Például:

\begin{lstlisting}[language=Velran]
let recent = catalog::queries::recent(db);
// tipus modulhataron at: models::Article
\end{lstlisting}

A \code{catalog/queries.vrn} tehát a \code{catalog::queries} névtérhez tartozik. A deklaráció saját modulján belül a rövid lokális név továbbra is használható. Így a helyi kód rövid marad, a forrásmodulok közötti dependency viszont látható.

A modulbetöltés továbbra is explicit és az application/package rooton belülre zárt: nincs \code{mod.vrn} fallback, automatikus könyvtár-discovery, \code{../}, \code{./} vagy wildcard import. A verified modulgráfon belül támogatottak a Rust-szerű \code{crate::}, \code{self::} és \code{super::} névtérgyökerek. A szűk, explicit \code{use path as alias;} import támogatott. A package-root re-export szándékosan szűkebb a Rustnál: \code{pub use path as alias;} csak már publikus modul-névteret tehet elérhetővé; wildcard, item-, private-module és private transitive-dependency re-export fail-closed módon tiltott.

A nyelvi műveletek -- például \code{text.trim()}, \code{x.sin()} és \code{values.len()} -- Rust-szerű metódusszintaxist használnak, amikor a művelet természetesen egy értékhez tartozik. Framework helper csak ott marad explicit függvény, ahol a Velran webes biztonsági szerződést ad hozzá, például \code{splitBounded(...)}. Ezek nem alkalmazásmodul-deklarációk, ezért továbbra is kvalifikáció nélkül használhatók.


\subsection*{Rust-szerű structok, inherent metódusok és konstruktorok}

Normál pure/domain számításhoz a Rust-alakú objektummodellt használd, ne egy második objektumszintaxist találj ki. A visibility explicit, és alapból private:

\begin{lstlisting}[language=Velran]
pub struct Summary {
    pub length: i64,
    hidden: i64,
}

impl Summary {
    pub fn new(length: i64) -> Self {
        return Self { length, hidden: 0 };
    }

    pub fn length_value(&self) -> i64 {
        return self.length;
    }
}

let summary = Summary::new(42);
let n = summary.length_value();
\end{lstlisting}

Az associated function, például \code{Type::new(...)}, illetve az inherent \code{impl}-en belüli \code{Self} ugyanazon verified pure/native lowering útvonalon halad, mint a normál függvény. A publikus metódus-aláírást a compiler rekurzívan visibility-ellenőrzi, ezért privát struct paraméteren vagy visszatérési típuson keresztül sem szivároghat ki, beleértve a beágyazott result/option alakokat is.

A jelenlegi biztonságos surface tudatosan szűkebb a teljes Rustnál: az általános mutable/owned receiver (például \code{\&mut self} vagy fogyasztó \code{self}), a trait impl és a generic impl továbbra is fail-closed, amíg a compiler az ownership/alias szemantikát a transaction, authorization, tenant isolation és idempotency határok gyengítése nélkül nem tudja bizonyítani.

\subsection*{Lokális package-ek és explicit dependencyk}

A workspace lokális, \code{velran.toml}-lal deklarált package-eket használhat. A dependency path-pinnelt, korlátozott és ciklusellenőrzött DAG-ként oldódik fel; a nyelvi package-rétegben nincs registry, Git dependency, build script vagy ambient hálózati feloldás. A közvetlen dependency saját package-névterén keresztül látható, a tranzitív dependency pedig privát marad a deklaráló package számára, kivéve a kifejezetten engedett publikus modul-névtér re-exportot.

\begin{lstlisting}
[package]
name = "demo_app"
kind = "app"
entry = "main.vrn"

[dependencies]
textkit = { path = "../libs/textkit" }
\end{lstlisting}

A library visibility és a webes authority két külön fogalom. A \code{pub} Rust-szerű library API-t tehet elérhetővé, de nem ad route/page/action/query/integration/filesystem/network jogosultságot. A library package-eket továbbra is a compiler-owned capability boundary védi, beleértve a fail-closed library web-authority szabályt.

\subsection*{Rust-szintaxis és framework-owned webes szintaxis}

Hasznos szabály: egy Rust-szerű compute nyelvet tanulj meg, és ehhez csak azokat a webspecifikus kiegészítéseket add hozzá, amelyeket a Velrannak kell birtokolnia. A határ szándékosan látható:

\begin{center}
\begin{tabularx}{\textwidth}{@{}L{0.30\textwidth}YY@{}}
\toprule
Terület & Rust-szerű forma & Velran-owned határ \\
\midrule
Függvény és lokális & \code{fn}, \code{let}, \code{let mut}, közvetlen assignment & \callable{page}, \callable{action}, \callable{query} \\
Érték és kollekció & \code{i64}, \code{f32}, \code{String}, \code{Vec<T>}, \code{Option<T>} & domain type, bounded request type és capability type \\
Control flow & \code{if}, \code{while}, exhaustív \code{match}, \code{?} & instruction/allocation budget és fail-closed verification \\
Web/data boundary & normál typed kifejezés & \code{route}, \code{html}, \code{sql}, \code{transaction}, validation/auth/cache policy \\
\bottomrule
\end{tabularx}
\end{center}

Ez a felosztás nem framework-megkerülés. Rust-szerű szintaxis csak akkor elfogadott, ha a natív lowering közben ugyanaz a verified, bounded webes szemantika megőrizhető.

\subsection*{Statement terminátorok}

A Velranban az egyszerű statementeket explicit pontosvessző zárja. A sortörés whitespace, nem statement-határ; nincs automatikus semicolon insertion. Ez determinisztikussá teszi a többsoros kifejezéseket és egyértelműbbé a generált, illetve auditált kódot.

\begin{lstlisting}[language=Velran]
let subtotal = price
    * quantity
    + shipping;
let mut retries = 0;
retries = retries + 1;
authorize article owner authorUsername or role Publisher;
return Ok(json(subtotal));
\end{lstlisting}

A blokkos statementeket és blokkos deklarációkat a kapcsos zárójel zárja, utánuk nincs pontosvessző:

\begin{lstlisting}[language=Velran]
let mut state = 0;
if ready {
    state = 1;
}
while state < 3 {
    state = state + 1;
}
\end{lstlisting}

Ugyanez érvényes tranzakción belül is: az önálló mutating query-hívás és a business-audit statement egyszerű statement, ezért \code{;} zárja. A top-level blokkos deklarációk -- például \code{model}, \callable{page}, \callable{action}, \callable{query} -- és a control-flow blokkok után nincs \code{;}. A blokk nélküli top-level deklarációk viszont pontosvesszősök: a \code{mod path;} és a \code{route ... => handler;} is explicit \code{;} terminátort kap. Többsoros route esetén is csak ez a záró terminátor fejezi be a deklarációt.

\section{Typed számítás webalkalmazás-kódban}

\subsection*{Kifejezések és aritmetika}

A Velran explicit operátorprecedenciát használ. Szorosabbtól lazább felé: unáris \code{!}; szorzás, osztás és maradékképzés; összeadás és kivonás; shiftek; összehasonlítások; bitenkénti AND/XOR/OR; logikai AND; logikai OR. Ha javítja az olvashatóságot, használj zárójelet.

\begin{lstlisting}[language=Velran]
let total = 12 + 3 * 4;
let remainder = 17 % 5;
let shifted = 3 << 2;
let masked = 12 & 10;
let allowed = ready && authorized;
let fallback = cached || fresh;
let denied = !allowed;
\end{lstlisting}

Az operátorszerződés szigorúan typed:

\begin{center}
\begin{tabularx}{\textwidth}{@{}l l Y@{}}
\toprule
\textbf{Operandusok} & \textbf{Operátorok} & \textbf{Eredmény} \\
\midrule
\code{i64}, azonos típus & \code{+ - * / \%} & checked \code{i64}. \\
\code{f32}, azonos típus & \texttt{+ - * / \%} & \code{f32}; véges eredmény szükséges. \\
\code{Decimal}, azonos típus & \code{+ - * / \%} & checked \code{Decimal}. \\
\code{i64}, azonos típus & \code{<< >> \& \^{} |} & shift/bitművelet, \code{i64}. \\
\code{bool}, azonos típus & \code{\&\& ||} & \code{bool}; short-circuit kiértékelés. \\
\code{bool} & \code{!} & logikai negáció. \\
\code{String}, azonos típus & \code{+} & budgetelt String-konkatenáció. \\
Numerikus, azonos típus & \code{< <= > >=} & \code{bool}. \\
Támogatott azonos skalártípus & \code{== !=} & \code{bool}. \\
\bottomrule
\end{tabularx}
\end{center}

Nincs implicit keverés \code{i64}, \code{f32} és \code{Decimal} között. Lebegőpontos számításhoz az integert explicit \code{value as f32} hívással alakítsd át. Overflow, nullával osztás/maradékképzés, hibás shift vagy nem véges \code{f32} eredmény fail-closed hibát ad.

A numerikus builtinok:

\begin{center}
\begin{tabularx}{\textwidth}{@{}l l Y@{}}
\toprule
\textbf{Függvény} & \textbf{Típus} & \textbf{Feladat} \\
\midrule
\code{x.sin()}, \code{x.cos()} & \code{f32 -> f32} & Trigonometrikus függvények. \\
\code{x.sqrt()} & \code{f32 -> f32} & Négyzetgyök; hibás tartomány elutasítva. \\
\code{x.abs()} & \code{i64 -> i64} vagy \code{f32 -> f32} & Abszolútérték. \\
\code{x.ln()}, \code{x.log10()} & \code{f32 -> f32} & Természetes/tízes alapú logaritmus. \\
\code{x.log(base)} & \code{f32, f32 -> f32} & Explicit alapú logaritmus. \\
\code{x.exp()} & \code{f32 -> f32} & Exponenciális függvény. \\
\code{x.powf(y)} & \code{f32, f32 -> f32} & Lebegőpontos hatványozás. \\
\code{x.round()/x.floor()/x.ceil()} & \code{f32 -> f32} & Kerekítési műveletek. \\
\code{x as f32} & \code{i64 -> f32} & Explicit integer--\code{f32} konverzió. \\
\code{monotonicNanos()} & \code{() -> i64} & Process-lokális monoton nanoszekundum-számláló. \\
\bottomrule
\end{tabularx}
\end{center}

Minden \code{f32} eredménynek végesnek kell maradnia. A \code{(-1.0f32).sqrt()} és \code{(-1.0f32).ln()} jellegű domainhibák NaN/végtelen helyett hibát adnak. A \code{monotonicNanos()} eltelt idő mérésére való, nem falióra-időbélyeg; a tőle függő public-cache kimenetet a compiler elutasítja.

\subsection*{Gyakori webes aritmetikai minták}

Webalkalmazásban ritkán kell mindenhol haladó matematika, néhány numerikus minta viszont folyamatosan visszatér. Ezeket typed értékekkel és látható kóddal kezeld, ne string-konverziókba vagy SQL-részletekbe rejtsd.

Lapozásnál validált integerekből számolj offsetet:

\begin{lstlisting}[language=Velran]
let pageSize = 25;
let offset = (page - 1) * pageSize;
\end{lstlisting}

Integer flag/mask esetén használhatók a bitműveletek, de üzleti feltételekhez az olvasható boolean logika az alap. A \code{&&} és \code{||} short-circuit módon működik:

\begin{lstlisting}[language=Velran]
let mayPublish = authenticated
    && ownsArticle
    && !locked;
\end{lstlisting}

Pénzhez \code{Decimal}-t használj. \code{f32}-t akkor válassz, ha maga a feladat lebegőpontos: például geometria, jelfeldolgozás, statisztika vagy explicit matematikai függvény. Pénzösszeget ne alakíts \code{f32}-re csak azért, hogy math builtint lehessen rá alkalmazni.

\subsection*{String- és szövegfüggvények}

A string műveletek typed, Unicode-tudatos builtinok a korlátozott runtime-on belül:

\begin{lstlisting}[language=Velran]
let cleaned = text.trim();
let left = text.trim_start();
let right = text.trim_end();
let normalized = cleaned.to_lowercase();
let chars = cleaned.chars().count();
let rewritten = normalized.replace(" ", "-");
let part = substring(cleaned, 1, 3);
let first = indexOf(cleaned, "vrn");
let last = lastIndexOf(cleaned, "a");
let ch = charAt("Velran", 1);
let repeated = "vrn".repeat(3);
let parts = splitBounded(rewritten, "-", 4096);
\end{lstlisting}

\begin{center}
\begin{tabularx}{\textwidth}{@{}Y Y@{}}
\toprule
\textbf{Függvény} & \textbf{Szignatúra} \\
\midrule
\code{text.chars().count()} & \code{String -> i64} \\
\code{text.trim()/trim\_start()/trim\_end()} & \code{String -> String} \\
\code{text.to\_lowercase()/to\_uppercase()} & \code{String -> String} \\
\code{text.contains(needle)} & \code{String, String -> bool} \\
\code{text.starts\_with(part) / text.ends\_with(part)} & \code{String, String -> bool} \\
\code{text.replace(search, replacement)} & \code{String, String, String -> String} \\
\code{splitBounded(text, delimiter, maxItems)} & \code{String, String, i64 -> Vec<String>} \\
\code{substring(text, start[, length])} & \code{String, i64[, i64] -> String} \\
\code{indexOf/lastIndexOf(text, needle)} & \code{String, String -> i64} \\
\code{charAt(text, index)} & \code{String, i64 -> String} \\
\code{text.repeat(count)} & \code{String, i64 -> String} \\
\bottomrule
\end{tabularx}
\end{center}

A \code{text.chars().count()}, \code{substring}, \code{indexOf}, \code{lastIndexOf} és \code{charAt} Unicode skalárérték-pozíciókat használ, nem UTF-8 byte indexet. Sikertelen keresésnél \code{-1} az eredmény; negatív vagy érvénytelen index fail-closed hibát ad. A \code{replace} nem enged üres keresőszöveget. A \code{splitBounded} nem enged üres delimitert, compile-time elemszám-korlátot kér a támogatott tartományban, és túllépéskor hibázik a csendes csonkolás helyett. A \code{repeat} az eredmény létrehozása előtt bekerül az allocation budgetbe. A \code{Vec<String>} typed compute collection marad \code{parts.len()} és ellenőrzött \code{parts[index]} használattal.

Korlátozott reguláris kifejezéses feldolgozáshoz elérhető a \code{regexMatch}, \code{regexReplace} és \code{regexCaptures}; ezek továbbra is a validációs és resource limitek alatt futnak.

\subsection*{Gyakori szövegkezelési minták}

A szövegmanipulációt kezeld boundary-feladatként. Keresési és filter inputot érdemes összehasonlítás előtt normalizálni, de az eredeti értéket tartsd meg, ha a domainben számít a pontos írásmód.

\begin{lstlisting}[language=Velran]
let queryText = q.trim();
let normalizedQuery = queryText.to_lowercase();
let hasVelran = normalizedQuery.contains("velran");
\end{lstlisting}

A \code{replace} determinisztikus helyettesítésre való, nem általános parser helyett:

\begin{lstlisting}[language=Velran]
let displayTag = tag.trim().replace("_", " ");
\end{lstlisting}

Szeletelésnél és pozíciókeresésnél a Velran Unicode scalar-value pozíciókat használ UTF-8 byte offset helyett. Így a \code{substring}, \code{charAt} és indexfüggvények alkalmazásszövegnél nem teszik ki a fejlesztőt encoding-részleteknek. Strukturált azonosítóknál -- URL, slug, email, dátum, UUID -- viszont a dedikált Velran típusokat és validátorokat használd, ne string függvényekből építsd újra a szabályokat.

\subsection*{Melyik eszközt válaszd alkalmazáskódban?}

\begin{center}
\begin{tabularx}{\textwidth}{@{}l Y@{}}
\toprule
\textbf{Webfejlesztési feladat} & \textbf{Preferált Velran megoldás} \\
\midrule
Keresési/filter szöveg normalizálása & \code{trim}, majd \code{lower}/\code{upper} csak akkor, ha az összehasonlítás case-insensitive. \\
Ismert token helyettesítése & \code{replace}; regexet csak valódi mintakereséshez használj. \\
Email, URL, slug, dátum, UUID & Dedikált typed érték és validáció, ne ad-hoc string parsing. \\
Pénz/összegzés & \code{Decimal}. \\
Lapozás/számlálók & \code{i64} checked aritmetikával. \\
Üzleti feltételek & \code{bool} short-circuit \code{&&}, \code{||}, \code{!} operátorokkal. \\
Tudományos/statisztikai matematika & \code{f32} és explicit math builtin; előbb döntsd el, hogy a lebegőpont valóban illik-e a domainhez. \\
Eltelt idő mérése & \code{monotonicNanos()}, soha ne falióra-időbélyegként. \\
\bottomrule
\end{tabularx}
\end{center}

\section{Domainértékek és handlerkód}

\subsection*{Modellek és üzleti típusok}

\begin{lstlisting}[language=Velran]
model Article {
    id: i64
    slug: Slug
    title: String
    body: String
}
\end{lstlisting}

A gyakori típusok: \code{String}, \code{i64}, \code{bool}, továbbá \code{Date}, \code{DateTime}, \code{Uuid}, \code{Decimal} és \code{Slug}. Pénzhez \code{Decimal}-t használj, ne lebegőpontos számot.

\subsection*{Page és action}

A page GET-szerű megjelenítési logika:

\begin{lstlisting}[language=Velran]
#[page]
fn show(ctx: PageContext, db: Db, id: i64)
    -> Result<Html, PageError> {
    let article = articleById(db, id)?;
    return Ok(html { <h1>{{ article.title }}</h1> });
}
\end{lstlisting}

Az action állapotváltoztató kérés kezelője:

\begin{lstlisting}[language=Velran]
#[action]
fn create(ctx: ActionContext, db: Db, title: String)
    -> Result<Redirect, PageError> {
    transaction db {
        insertArticle(tx, title)?;
    }
    return Ok(redirect(adminArticles()));
}
\end{lstlisting}

A példa egy \code{adminArticles} nevű GET route-ot feltételez; a redirect cél route call, nem kézzel másolt URL string.

\subsection*{Változók és hibaterjedés}

A \code{let} immutable typed lokális értéket hoz létre. Ha az értéket módosítani kell, Rust-szerű \code{let mut} deklarációt használj, majd normál assignmentet, például \code{count = count + 1} vagy \code{samples[i] = value}. A statikus típus nem változhat, a bounds check és az erőforrás-elszámolás megmarad, a legacy \code{set} szintaxis pedig tiltott. A \code{?} a typed \code{Result} hibáját továbbítja.

\subsection*{Enum és exhaustive match}

Zárt választási halmazhoz ne szabad szöveget használj, hanem enumot. Ez form-, route-, JSON- és domain-contractként is erősebb. Az enum egyben zárt sum type: ha a vezérlés varianttól függ, használj exhaustive \code{match}-et.

\begin{lstlisting}[language=Velran]
enum PublishState {
    Draft
    Review
    Published
}

match state {
    Draft => { return Ok(json("draft")); }
    Review => { return Ok(json("review")); }
    Published => { return Ok(json("published")); }
}
\end{lstlisting}

Ennél a security-releváns formánál nincs wildcard fallthrough. Ha új variant kerül a típusba, a régi match-ek addig nem fordulnak le, amíg az új állapot nincs explicit kezelve. A tranzakciós fejezet ugyanezt használja a \code{CommitUnknown} esetén.

\subsection*{Object}

Összetartozó model/query/page műveletek domain object alatt csoportosíthatók. Ez különösen nagyobb alkalmazásban segít a domain API körülhatárolásában.


\section{HTTP- és prezentációs határ}

\subsection*{Route}

\begin{lstlisting}[language=Velran]
route article GET "/cikk/:slug<Slug>" public => article;
\end{lstlisting}

A path paraméter típusa a route-ban deklarált. Hibás input nem jut el a handlerig typed értékként.

\subsection*{HTML escape}

\begin{lstlisting}[language=Velran]
return Ok(html {
    <h1>{{ article.title }}</h1>
});
\end{lstlisting}

A \code{{ ... }} HTML-escaped. Adatbázisból érkező string sem kivétel. Raw HTML API nincs.

\subsection*{Lista és optional érték}

A listát iterálni kell, az optional modelt pedig explicit jelenlétvizsgálat után lehet mezőzni. A compiler nem engedi, hogy \code{Vec<Model>} vagy ellenőrizetlen \code{Option<Model>} úgy viselkedjen, mintha egyetlen rekord lenne.

\begin{lstlisting}[language=Velran]
@for article in articles {
    <li>{{ article.title }}</li>
}

@if article {
    <h1>{{ article.title }}</h1>
}
\end{lstlisting}

A query-cardinality pontos futásidejű szerződését -- például hogy a \code{Option<Model>} több sor esetén hibát jelent -- az adatbázis-fejezet tárgyalja.

\subsection*{Typed URL}

\begin{lstlisting}[language=Velran]
<a @href(article, article.slug)>Megnyitas</a>
<form method="post" @action(deleteArticle, article.id)>
\end{lstlisting}

\subsection*{Preferáld a route helpereket}

Saját alkalmazás-route-ok URL-jét ne másold kézzel több helyre. HTML-ben POST formhoz az \code{@action(...)} és linkhez az \code{@href(...)} a preferált minta, mert a route neve és typed paraméterei maradnak a kapcsolat forrásai. A beépített runtime endpointok -- például a \code{/__velran/auth/logout} -- természetesen nem alkalmazás-route helperek.


Ne építs dinamikus \code{href="{{ value }}"} linket. A route helper ismeri a célt és a paramétertípust.

\subsection*{Component és layout}

\begin{lstlisting}[language=Velran]
#[layout]
fn Main(title: String) -> Html {
    html {
        <html>
            <head><title>{{ title }}</title></head>
            <body><main>@content</main></body>
        </html>
    }
}
\end{lstlisting}

Page-ben:

\begin{lstlisting}[language=Velran]
return Ok(html {
    @layout(Main, article.title) {
        <article><h1>{{ article.title }}</h1></article>
    }
});
\end{lstlisting}

A layout pontosan egy \code{@content} slotot tartalmazhat; component/layout csak explicit typed paramétereket lát.

\section{Gyakorlat: döntsd el, kié a szintaxis}
\paragraph{Gyakorlási mód: vezetett.} Használd az előszóban meghatározott vezetett módszert.

Nézz meg egy page-et, actiont, query-t, route-ot, komponenst és egy hétköznapi kifejezést. Mindegyiknél döntsd el, hogy normál Rust-szerű szintaxisról vagy Velran framework által birtokolt webes határról van-e szó. Ez segít új szintaxis kitalálása helyett helyesen következtetni: ahol lehet, Rust-szerű alakot várj; külön forma főleg security-, HTTP-, resource- vagy framework-contractnál indokolt.

\section{Tipikus hiba: kényelmi szintaxist találunk ki}
Ne feltételezz dokumentálatlan aliasokat, legacy callable alakokat vagy egyedi collection-típusneveket. A nyelv szándékosan közelít a Rust szintaxisához, és csak azokat a framework konstrukciókat tartja meg, amelyek a webes határokat explicitté és ellenőrizhetővé teszik.

\section{Folyamatos mintaprojekt: Secure Notes}
Definiáld a jegyzetlista page-et és a létrehozó actiont az aktuális callable attribútumokkal, típusos route inputokkal és explicit access policy-val. Az üzleti adat maradjon típusos érték, ne handlerek között továbbadott strukturálatlan string.
